\FigureLayoutDeclare{img/2015/national_paper_2_liberal/q6_fig1.png}{3.0cm}{13bp 8bp 21bp 12bp}

\chapter{2015年全国II卷（文）}

\section{单选题}

\begin{problem}
已知集合 \(A = \{x \mid -1 < x < 2\}\)，\(B = \{x \mid 0 < x < 3\}\)，则 \(A \cup B =\)（\quad）

\choices
    {\((-1,3)\)}
    {\((-1,0)\)}
    {\((0, 2)\)}
    {\((2,3)\)}
\end{problem}

\answer{A}

\begin{solution}
由 \(A=(-1,2)\)，\(B=(0,3)\)，两区间有公共部分，且左端点为 \(-1\)，右端点为 \(3\)，所以
\[
A\cup B=(-1,3)
\]
故选 A.
\end{solution}

\begin{problem}
若 \(a\) 为实数，且 \(\frac{2 + a\i}{1 + \i} = 3 + \i\)，则 \(a =\)（\quad）

\choices
    {\(-4\)}
    {\(-3\)}
    {3}
    {4}
\end{problem}

\answer{D}

\begin{solution}
由复数等式得
\[
2+a\i=(3+\i)(1+\i)=2+4\i
\]
比较虚部，得 \(a=4\). 故选 D.
\end{solution}

\begin{problem}
根据下面给出的 \(2004\) 年至 \(2013\) 年我国二氧化硫年排放量（单位：万吨）柱形图，以下结论中不正确的是（\quad）

\texfigure[max width=0.68\linewidth]{img_repaint/2015/national_paper_2_liberal/q3_fig1.tex}

\choices
    {逐年比较，\(2008\) 年减少二氧化硫排放量的效果最显著}
    {\(2007\) 年我国治理二氧化硫排放显现成效}
    {\(2006\) 年以来我国二氧化硫年排放量呈减少趋势}
    {\(2006\) 年以来我国二氧化硫年排放量与年份正相关}
\end{problem}

\answer{D}

\begin{solution}
由柱形图可见，\(2006\) 年以后排放量总体下降，虽然 \(2011\) 年较 \(2010\) 年略有上升；逐年比较相邻年份的下降量，\(2008\) 年的下降最明显. 因此 A、B、C 均可由图得到，而“排放量与年份正相关”与总体下降趋势相矛盾. 故选 D.
\end{solution}

\begin{problem}
已知 \(\bs{a} = (1, -1)\)，\(\bs{b} = (-1, 2)\)，则 \((2\bs{a} + \bs{b}) \cdot \bs{a} =\)（\quad）

\choices
    {\(-1\)}
    {0}
    {1}
    {2}
\end{problem}

\answer{C}

\begin{solution}
\[
2\bs{a}+\bs{b}=2(1,-1)+(-1,2)=(1,0)
\]
所以
\[
(2\bs{a}+\bs{b})\cdot\bs{a}=(1,0)\cdot(1,-1)=1
\]
故选 C.
\end{solution}

\begin{problem}
设 \(S_{n}\) 是等差数列 \(\{a_{n}\}\) 的前 \(n\) 项和. 若 \(a_{1} + a_{3} + a_{5} = 3\)，则 \(S_{5} =\)（\quad）

\choices
    {5}
    {7}
    {9}
    {11}
\end{problem}

\answer{A}

\begin{solution}
等差数列中 \(a_{1}+a_{5}=2a_{3}\)，所以
\[
a_{1}+a_{3}+a_{5}=3a_{3}=3
\]
从而 \(a_{3}=1\). 又等差数列前五项的平均数为中间项 \(a_{3}\)，故
\[
S_{5}=5a_{3}=5
\]
故选 A.
\end{solution}

\begin{problem}
一个正方体被一个平面截去一部分后，剩余部分的三视图如图，则截去部分体积与剩余部分体积的比值为（\quad）

\bitmapfigure[max width=0.36\linewidth]{img/2015/national_paper_2_liberal/q6_fig1.png}

\choices
    {\(\frac{1}{8}\)}
    {\(\frac{1}{7}\)}
    {\(\frac{1}{6}\)}
    {\(\frac{1}{5}\)}
\end{problem}

\answer{D}

\begin{solution}
设正方体棱长为 \(a\). 由三个三视图中的对角线可知，截去部分是以正方体一个顶点为直角顶点、三条互相垂直的棱长均为 \(a\) 的三棱锥. 以其中两条棱为底边，底面积为 \(a^{2}/2\)，第三条棱为高，因此
\[
V_{\text{截去}}=\frac{1}{3}\cdot\frac{a^{2}}{2}\cdot a=\frac{a^{3}}{6}
\]
剩余部分体积为
\[
V_{\text{剩余}}=a^{3}-\frac{a^{3}}{6}=\frac{5a^{3}}{6}
\]
所以体积比为
\[
\frac{V_{\text{截去}}}{V_{\text{剩余}}}=\frac{a^{3}/6}{5a^{3}/6}=\frac{1}{5}
\]
故选 D.
\end{solution}

\begin{problem}
已知三点 \(A(1,0)\)，\(B(0,\sqrt{3})\)，\(C(2,\sqrt{3})\)，则 \(\triangle ABC\) 外接圆的圆心到原点的距离为（\quad）

\choices
    {\(\frac{5}{3}\)}
    {\(\frac{\sqrt{21}}{3}\)}
    {\(\frac{2\sqrt{5}}{3}\)}
    {\(\frac{4}{3}\)}
\end{problem}

\answer{B}

\begin{solution}
因为 \(BC\) 是水平线段，且其中点为 \((1,\sqrt{3})\)，所以三角形外接圆圆心在直线 \(x=1\) 上. 设圆心为 \(P(1,p)\)，由 \(PA=PB\) 得
\[
p^{2}=1+(p-\sqrt{3})^{2}
\]
解得 \(p=\dfrac{2\sqrt{3}}{3}\)，故
\[
OP=\sqrt{1+\left(\frac{2\sqrt{3}}{3}\right)^{2}}=\frac{\sqrt{21}}{3}
\]
故选 B.
\end{solution}

\begin{problem}
下边程序框图的算法思路源于我国古代数学名著《九章算术》中的“更相减损术”. 执行该程序框图，若输入的 \(a\)，\(b\) 分别为 \(14\)，\(18\)，则输出的 \(a =\)（\quad）

\texfigure[max width=0.42\linewidth]{img_repaint/2015/national_paper_2_liberal/q8_fig1.tex}

\choices
    {0}
    {2}
    {4}
    {14}
\end{problem}

\answer{B}

\begin{solution}
程序每次在 \(a\ne b\) 时用较大数减去较小数. 初始为 \((a,b)=(14,18)\)，依次得到
\[
(14,18)\to(14,4)\to(10,4)\to(6,4)\to(2,4)\to(2,2)
\]
此时 \(a=b\)，程序输出 \(a=2\). 故选 B.
\end{solution}

\begin{problem}
已知等比数列 \(\{a_{n}\}\) 满足 \(a_{1} = \frac{1}{4}\)，\(a_{3}a_{5} = 4(a_{4} - 1)\)，则 \(a_{2} =\)（\quad）

\choices
    {2}
    {1}
    {\(\frac{1}{2}\)}
    {\(\frac{1}{8}\)}
\end{problem}

\answer{C}

\begin{solution}
设等比数列的公比为 \(q\). 则
\(a_{3}=\frac{q^{2}}{4}\)，\(a_{4}=\frac{q^{3}}{4}\)，\(a_{5}=\frac{q^{4}}{4}\).
代入条件，得
\[
\frac{q^{6}}{16}=4\left(\frac{q^{3}}{4}-1\right)=q^{3}-4
\]
即
\[
(q^{3}-8)^{2}=0
\]
所以 \(q=2\)，从而 \(a_{2}=a_{1}q=\dfrac12\). 故选 C.
\end{solution}

\begin{problem}
已知 \(A\)，\(B\) 是球 \(O\) 的球面上两点，\(\angle AOB = 90^{\circ}\)，\(C\) 为该球面上的动点，若三棱锥 \(O-ABC\) 体积的最大值为 \(36\)，则球 \(O\) 的表面积为（\quad）

\choices
    {\(36\pi\)}
    {\(64\pi\)}
    {\(144\pi\)}
    {\(256\pi\)}
\end{problem}

\answer{C}

\begin{solution}
设球 \(O\) 的半径为 \(R\). 以三角形 \(OAB\) 为底面、点 \(C\) 为顶点，则
\[
S_{\triangle OAB}=\frac12 OA\cdot OB\sin\angle AOB=\frac{R^{2}}{2}
\]
平面 \(OAB\) 经过球心，球面上点 \(C\) 到该平面的距离最大为 \(R\)，因此三棱锥体积的最大值为
\[
V_{\max}=\frac13\cdot\frac{R^{2}}{2}\cdot R=\frac{R^{3}}{6}
\]
由 \(V_{\max}=36\) 得 \(R=6\)，所以球的表面积为
\[
4\pi R^{2}=144\pi
\]
故选 C.
\end{solution}

\begin{problem}
如图，长方形 \(ABCD\) 的边 \(AB = 2\)，\(BC = 1\)，\(O\) 是 \(AB\) 的中点，点 \(P\) 沿着边 \(BC\)，\(CD\) 与 \(DA\) 运动，记 \(\angle BOP = x\). 将动点 \(P\) 到 \(A\)，\(B\) 两点距离之和表示为 \(x\) 的函数 \(f(x)\)，则 \(y = f(x)\) 的图象大致为（\quad）

\bitmapfigure[max width=0.42\linewidth]{img/2015/national_paper_2_liberal/q11_fig1.png}

\choices
    {\choicebitmap[width=0.15\paperwidth]{img/2015/national_paper_2_liberal/q11_opt_a.png}}
    {\choicebitmap[width=0.15\paperwidth]{img/2015/national_paper_2_liberal/q11_opt_b.png}}
    {\choicebitmap[width=0.15\paperwidth]{img/2015/national_paper_2_liberal/q11_opt_c.png}}
    {\choicebitmap[width=0.15\paperwidth]{img/2015/national_paper_2_liberal/q11_opt_d.png}}
\end{problem}

\answer{B}

\begin{solution}
取 \(A=(0,0)\)，\(B=(2,0)\)，\(C=(2,1)\)，\(D=(0,1)\)，则 \(O=(1,0)\). 当 \(P\) 沿 \(BC\) 运动时，令 \(P=(2,t)\)，有 \(0\leqslant t\leqslant1\)，且 \(t=\tan x\)，于是
\(f(x)=\sqrt{4+\tan^{2}x}+\tan x\)，\(0\leqslant x\leqslant\frac{\pi}{4}\)
该段严格递增.

当 \(P\) 沿 \(CD\) 运动时，令 \(P=(u,1)\)，则 \(x\) 从 \(\dfrac{\pi}{4}\) 增加到 \(\dfrac{3\pi}{4}\)，且
\[
f(x)=\sqrt{u^{2}+1}+\sqrt{(u-2)^{2}+1}
\]
记右端为 \(F(u)\)，则
\[
F''(u)=\frac{1}{(u^{2}+1)^{3/2}}+\frac{1}{((u-2)^{2}+1)^{3/2}}>0
\]
所以 \(F(u)\) 严格凸，且 \(F(2-u)=F(u)\)，从而在 \(u=1\) 时取得最小值 \(2\sqrt2\). 点 \(P=(u,1)\) 关于直线 \(x=1\) 对称时，\(\angle BOP\) 分别为 \(x\) 与 \(\pi-x\)，故这一段图象关于 \(x=\dfrac\pi2\) 对称，先下降后上升.

当 \(P\) 沿 \(DA\) 运动时，\(x\) 从 \(\dfrac{3\pi}{4}\) 增加到 \(\pi\)，点 \(P\) 逐渐靠近 \(A\)，故距离和递减，且在 \(A\) 处为 \(2\). 因而图象从 \(x=0\) 处的 \(2\) 上升，在 \(\dfrac\pi2\) 附近下降后再上升，最后下降到 \(x=\pi\) 处的 \(2\)，符合选项 B. 故选 B.
\end{solution}

\begin{problem}
设函数 \( f(x) = \ln (1 + |x|) - \frac{1}{1 + x^2} \)，则使得 \( f(x) > f(2x - 1) \) 成立的 \(x\) 的取值范围是（\quad）

\choices
    {\(\left(\frac{1}{3}, 1\right)\)}
    {\(\left(-\infty, \frac{1}{3}\right) \cup (1, +\infty)\)}
    {\(\left(-\frac{1}{3}, \frac{1}{3}\right)\)}
    {\(\left(-\infty, -\frac{1}{3}\right) \cup \left(\frac{1}{3}, +\infty\right)\)}
\end{problem}

\answer{A}

\begin{solution}
令
\(g(t)=\ln(1+t)-\frac{1}{1+t^{2}}\)，\(t\geqslant0\).
则
\[
g'(t)=\frac{1}{1+t}+\frac{2t}{(1+t^{2})^{2}}>0
\]
所以 \(g(t)\) 在 \([0,+\infty)\) 上严格递增. 又 \(f(x)=g(|x|)\)，故
\[
f(x)>f(2x-1)\Longleftrightarrow |x|>|2x-1|
\]
两边平方并整理，得
\[
x^{2}>(2x-1)^{2}\Longleftrightarrow 3x^{2}-4x+1<0
\]
该二次式的两个根为 \(\dfrac13\) 和 \(1\)，所以
\[
\frac13<x<1
\]
故选 A.
\end{solution}

\section{填空题}

\begin{problem}
已知函数 \( f(x) = ax^3 - 2x \) 的图象过点 \((-1,4)\)，则 \(a =\)\fillinblank{}.
\end{problem}

\begin{answer}
-2
\end{answer}

\begin{solution}
将点 \((-1,4)\) 代入函数方程，得
\[
4=a(-1)^{3}-2(-1)=-a+2
\]
所以 \(a=-2\).
\end{solution}

\begin{problem}
若 \(x\)，\(y\) 满足约束条件 \(\begin{cases}
x + y - 5 \leqslant 0,\\
2x - y - 1 \geqslant 0,\\
x - 2y + 1 \leqslant 0,
\end{cases}\) 则 \(z = 2x + y\) 的最大值为\fillinblank{}.
\end{problem}

\begin{answer}
8
\end{answer}

\begin{solution}
三个边界直线分别为
\(y=5-x\)，\(y=2x-1\)，\(y=\frac{x+1}{2}\).
它们两两相交得到可行域的三个顶点
\((2,3)\)，\((3,2)\)，\((1,1)\).
在这三个顶点处，目标函数 \(z=2x+y\) 的值分别为 \(7\)，\(8\)，\(3\). 线性目标函数在三角形可行域上取最大值于顶点，故最大值为 \(8\).
\end{solution}

\begin{problem}
已知双曲线过点 \((4, \sqrt{3})\)，且渐近线方程为 \(y = \pm \frac{1}{2}x\)，则该双曲线的标准方程为\fillinblank{}.
\end{problem}

\begin{answer}
\(\dfrac{x^{2}}{4}-y^{2}=1\)
\end{answer}

\begin{solution}
若双曲线为横轴型 \(\dfrac{x^{2}}{a^{2}}-\dfrac{y^{2}}{b^{2}}=1\)，其渐近线斜率为 \(\pm\dfrac{b}{a}\)，故 \(b=\dfrac{a}{2}\). 代入点 \((4,\sqrt{3})\)，得
\[
\frac{16}{a^{2}}-\frac{3}{a^{2}/4}=1
\]
从而 \(a^{2}=4\)，\(b^{2}=1\).

若为纵轴型，则方程为 \(\dfrac{y^{2}}{a^{2}}-\dfrac{x^{2}}{b^{2}}=1\)，且由渐近线斜率有 \(a^{2}/b^{2}=1/4\). 代入点 \((4,\sqrt{3})\) 后左端为
\[
\frac{3}{a^{2}}-\frac{16}{b^{2}}=\frac{3}{a^{2}}-\frac{4}{a^{2}}=-\frac{1}{a^{2}}<0
\]
不可能等于 \(1\). 因而双曲线的方程为
\[
\frac{x^{2}}{4}-y^{2}=1
\]
\end{solution}

\begin{problem}
已知曲线 \( y = x + \ln x \) 在点 \((1, 1)\) 处的切线与曲线 \( y = ax^2 + (a + 2)x + 1 \) 相切，则 \(a =\)\fillinblank{}.
\end{problem}

\begin{answer}
8
\end{answer}

\begin{solution}
曲线 \(y=x+\ln x\) 在 \(x=1\) 处的导数为 \(1+\dfrac11=2\)，所以切线为
\(y-1=2(x-1)\)，\(y=2x-1\).
令该直线与第二条曲线相交，得
\[
ax^{2}+(a+2)x+1=2x-1
\]
即
\[
ax^{2}+ax+2=0
\]
相切要求有重根，且 \(a=0\) 时方程退化为 \(2=0\)，没有交点，因此 \(a\ne0\). 由判别式为零，得
\[
a^{2}-8a=0\Longrightarrow a=8
\]
\end{solution}

\section{解答题}

\begin{problem}
\(\triangle ABC\) 中，\(D\) 是 \(BC\) 上的点，\(AD\) 平分 \(\angle BAC\)，\(BD = 2DC\).
\begin{enumerate}
\item 求 \(\frac{\sin\angle B}{\sin\angle C}\)；
\item 若 \(\angle BAC = 60^{\circ}\)，求 \(\angle B\).
\end{enumerate}
\end{problem}

\begin{solution}
\begin{enumerate}
\item 由角平分线定理，
\[
\frac{AB}{AC}=\frac{BD}{DC}=2
\]
在 \(\triangle ABC\) 中由正弦定理，\(\dfrac{AB}{AC}=\dfrac{\sin\angle C}{\sin\angle B}\)，所以
\[
\frac{\sin\angle B}{\sin\angle C}=\frac12
\]

\item 设 \(\angle B=\beta\)，则 \(\angle C=120^{\circ}-\beta\). 由第（1）问，
\[
\frac{\sin\beta}{\sin(120^{\circ}-\beta)}=\frac12
\]
利用正弦差角公式，得
\[
2\sin\beta=\sin(120^{\circ}-\beta)=\frac{\sqrt3}{2}\cos\beta+\frac12\sin\beta
\]
从而 \(3\sin\beta=\sqrt3\cos\beta\)，即 \(\tan\beta=\dfrac{\sqrt3}{3}\). 因为三角形内角满足 \(0^{\circ}<\beta<120^{\circ}\)，故 \(\beta=30^{\circ}\)，所以 \(\angle B=30^{\circ}\).
\end{enumerate}
\end{solution}

\begin{problem}
某公司为了了解用户对其产品的满意度，从 \(A\)，\(B\) 两地区分别随机调查了 \(40\) 个用户，根据用户对其产品的满意度的评分，得到 \(A\) 地区用户满意度评分的频率分布直方图和 \(B\) 地区用户满意度评分的频率分布表.

\begin{center}
\(B\) 地区用户满意度评分的频率分布表
\begin{tabular}{c|ccccc}
\hline
满意度评分分组 & \([50, 60)\) & \([60, 70)\) & \([70, 80)\) & \([80, 90)\) & \([90, 100]\) \\
\hline
频数 & 2 & 8 & 14 & 10 & 6 \\
\hline
\end{tabular}
\end{center}

\begin{enumerate}
\item 在图中作出 \(B\) 地区用户满意度评分的频率分布直方图，并通过此图比较两地区满意度评分的平均值及分散程度；
\item 根据用户满意度评分，将用户的满意度评分分为三个等级：

\begin{center}
\begin{tabular}{c|ccc}
\hline
满意度评分 & 低于 \(70\) 分 & \(70\) 分到 \(89\) 分 & 不低于 \(90\) 分 \\
\hline
满意度等级 & 不满意 & 满意 & 非常满意 \\
\hline
\end{tabular}
\end{center}

估计哪个地区的用户的满意度等级为不满意的概率大，说明理由.
\end{enumerate}

\texfigure[max width=0.62\linewidth]{img_repaint/2015/national_paper_2_liberal/q18_fig1.tex}
\end{problem}

\begin{solution}
\begin{enumerate}
\item 每组组距均为 \(10\). 由频数除以样本数得到 \(B\) 地区各组频率，再除以组距，频率分布直方图各柱高依次为
\(0.005\)，\(0.020\)，\(0.035\)，\(0.025\)，\(0.015\)
对应的评分区间依次为 \([50,60)\)，\([60,70)\)，\([70,80)\)，\([80,90)\)，\([90,100]\).

用各组组中值估计平均分，\(A\) 地区的平均分为
\[
45\times0.10+55\times0.20+65\times0.30+75\times0.20+85\times0.15+95\times0.05=67.5
\]
\(B\) 地区的平均分为
\[
55\times0.05+65\times0.20+75\times0.35+85\times0.25+95\times0.15=77.5
\]
所以 \(B\) 地区的平均满意度评分较高. 以组中值估计方差，两个地区分别为
\[
\begin{aligned}
s_A^2&=0.10(45-67.5)^2+0.20(55-67.5)^2+0.30(65-67.5)^2\\
&\quad+0.20(75-67.5)^2+0.15(85-67.5)^2+0.05(95-67.5)^2=178.75,\\
s_B^2&=0.05(55-77.5)^2+0.20(65-77.5)^2+0.35(75-77.5)^2\\
&\quad+0.25(85-77.5)^2+0.15(95-77.5)^2=118.75.
\end{aligned}
\]
因此 \(s_B^2<s_A^2\)，\(B\) 地区评分的分散程度较小.

\item \(A\) 地区评分低于 \(70\) 分的频率估计为
\[
0.10+0.20+0.30=0.60
\]
\(B\) 地区评分低于 \(70\) 分的频率为
\[
\frac{2+8}{40}=0.25
\]
因为 \(0.60>0.25\)，所以估计 \(A\) 地区用户的满意度等级为不满意的概率较大.
\end{enumerate}
\end{solution}

\begin{problem}
如图，长方体 \(ABCD - A_{1}B_{1}C_{1}D_{1}\) 中，\(AB = 16\)，\(BC = 10\)，\(AA_{1} = 8\)，点 \(E\)，\(F\) 分别在 \(A_{1}B_{1}\)，\(D_{1}C_{1}\) 上，\(A_{1}E = D_{1}F = 4\). 过点 \(E\)，\(F\) 的平面 \(\alpha\) 与此长方体的面相交，交线围成一个正方形.

\begin{enumerate}
\item 在图中画出这个正方形；
\item 求平面 \(\alpha\) 把该长方体分成的两部分体积的比值.
\end{enumerate}

\bitmapfigure[max width=0.66\linewidth]{img/2015/national_paper_2_liberal/q19_fig1.png}
\end{problem}

\begin{solution}
在 \(AB\) 上取点 \(K\)，在 \(DC\) 上取点 \(L\)，使 \(AK=DL=10\)，连接 \(EK\)，\(FL\)，则交线围成的正方形为 \(EFKL\). 事实上，\(EF=KL=BC=10\)，且在侧面 \(ABB_{1}A_{1}\) 中
\[
EK=\sqrt{(10-4)^{2}+8^{2}}=10
\]
同理 \(FL=10\)，并且 \(EF\) 与 \(EK\) 垂直，所以该截面为边长 \(10\) 的正方形.

取坐标系使 \(A=(0,0,0)\)，\(B=(16,0,0)\)，\(D=(0,10,0)\)，\(A_{1}=(0,0,8)\). 则
\(E=(4,0,8)\)，\(F=(4,10,8)\)，\(K=(10,0,0)\)，\(L=(10,10,0)\).
平面 \(\alpha\) 的方程为 \(4x+3z=40\). 其下方一侧的体积为
\[
V_{1}=4\times10\times8+\int_{4}^{10}10\cdot\frac{40-4x}{3}\,\mathrm{d}x=320+240=560
\]
长方体体积为 \(16\times10\times8=1280\)，另一侧体积为 \(V_{2}=1280-560=720\). 因此两部分体积之比为
\[
V_{1}:V_{2}=560:720=7:9
\]
\end{solution}

\begin{problem}
已知椭圆 \(C: \frac{x^2}{a^2} + \frac{y^2}{b^2} = 1\)（\(a > b > 0\)） 的离心率为 \(\frac{\sqrt{2}}{2}\)，点 \((2, \sqrt{2})\) 在 \(C\) 上.
\begin{enumerate}
\item 求 \(C\) 的方程；
\item 直线 \(l\) 不经过原点 \(O\)，且不平行于坐标轴，\(l\) 与 \(C\) 有两个交点 \(A\)，\(B\)，线段 \(AB\) 中点为 \(M\)，证明：直线 \(OM\) 的斜率与直线 \(l\) 的斜率乘积为定值.
\end{enumerate}
\end{problem}

\begin{solution}
\begin{enumerate}
\item 设椭圆半长轴、半短轴和半焦距分别为 \(a\)，\(b\)，\(c\). 由离心率得
\(\frac{c}{a}=\frac{\sqrt2}{2}\)，\(b^{2}=a^{2}-c^{2}=\frac{a^{2}}{2}\).
又因为 \((2,\sqrt2)\) 在椭圆上，
\[
\frac{4}{a^{2}}+\frac{2}{b^{2}}=1
\]
代入 \(b^{2}=a^{2}/2\)，得 \(a^{2}=8\)，\(b^{2}=4\). 所以
\[
C:\frac{x^{2}}{8}+\frac{y^{2}}{4}=1
\]

\item 设直线 \(l\) 的方程为 \(y=kx+m\). 因为 \(l\) 不平行于坐标轴且不过原点，所以 \(k\ne0\)，\(m\ne0\). 设 \(A(x_{1},y_{1})\)，\(B(x_{2},y_{2})\)，则 \(x_{1}\)，\(x_{2}\) 是方程
\[
x^{2}+2(kx+m)^{2}=8
\]
的两个根，即
\[
(1+2k^{2})x^{2}+4kmx+2m^{2}-8=0
\]
由韦达定理及中点坐标公式，
\[
x_{M}=\frac{x_{1}+x_{2}}{2}=-\frac{2km}{1+2k^{2}}
\]
且
\[
y_{M}=kx_{M}+m=\frac{m}{1+2k^{2}}
\]
因此
\[
k_{OM}=\frac{y_{M}}{x_{M}}=-\frac{1}{2k}
\]
从而直线 \(OM\) 与直线 \(l\) 的斜率乘积为
\[
k_{OM}\cdot k=-\frac12
\]
故该定值为 \(-\dfrac12\).
\end{enumerate}
\end{solution}

\begin{problem}
已知 \(f(x) = \ln x + a(1 - x)\).
\begin{enumerate}
\item 讨论 \(f(x)\) 的单调性；
\item 当 \(f(x)\) 有最大值，且最大值大于 \(2a - 2\) 时，求 \(a\) 的取值范围.
\end{enumerate}
\end{problem}

\begin{solution}
\begin{enumerate}
\item 定义域为 \((0,+\infty)\)，且
\[
f'(x)=\frac1x-a=\frac{1-ax}{x}
\]
当 \(a\leqslant0\) 时，\(f'(x)>0\)，所以 \(f(x)\) 在 \((0,+\infty)\) 上单调递增. 当 \(a>0\) 时，\(f'(x)>0\) 在 \(\left(0,\dfrac1a\right)\) 上成立，\(f'(x)<0\) 在 \(\left(\dfrac1a,+\infty\right)\) 上成立，所以 \(f(x)\) 先增后减.

\item 当 \(a\leqslant0\) 时，\(f(x)\) 在无穷远处趋于 \(+\infty\)，没有最大值. 因此必须有 \(a>0\)，此时最大值为
\[
f\left(\frac1a\right)=-\ln a+a-1
\]
条件给出
\[
-\ln a+a-1>2a-2\Longleftrightarrow \ln a+a<1
\]
函数 \(h(a)=\ln a+a\) 在 \((0,+\infty)\) 上严格递增，且 \(h(1)=1\)，所以
\[
0<a<1
\]
\end{enumerate}
\end{solution}

\section{选考题：解答题}

\begin{problem}
如图，\(O\) 为等腰三角形 \(ABC\) 内一点，\(\odot O\) 与 \(\triangle ABC\) 的底边 \(BC\) 交于 \(M\)，\(N\) 两点，与底边上的高 \(AD\) 交于点 \(G\)，且与 \(AB\)，\(AC\) 分别相切于 \(E\)，\(F\) 两点.

\begin{enumerate}
\item 证明：\(EF \parallel BC\)；
\item 若 \(AG\) 等于 \(\odot O\) 的半径，且 \(AE = MN = 2\sqrt{3}\)，求四边形 \(EBCF\) 的面积.
\end{enumerate}

\bitmapfigure[max width=0.38\linewidth]{img/2015/national_paper_2_liberal/q22_fig1.png}
\end{problem}

\begin{solution}
\begin{enumerate}
\item 因为 \(ABC\) 是等腰三角形，\(AD\) 是顶角平分线，也是对称轴. 圆与 \(AB\)，\(AC\) 分别相切于 \(E\)，\(F\)，故圆心 \(O\) 到两腰的距离相等，且 \(O\) 在 \(AD\) 上. 于是直角三角形 \(AOE\) 与 \(AOF\) 全等，从而 \(AE=AF\). 又 \(AB=AC\)，所以
\[
\frac{AE}{AB}=\frac{AF}{AC}
\]
由三角形相似的逆定理，得 \(EF\parallel BC\).

\item 设圆的半径为 \(r\). 由图中 \(G\) 为 \(AD\) 上圆的上交点，且 \(AG=r\)，又 \(OG=r\)，所以
\[
AO=AG+GO=2r
\]
由于 \(AE\) 是从点 \(A\) 引圆的切线，\(OE\perp AB\)，在直角三角形 \(AOE\) 中
\[
AE^{2}=AO^{2}-OE^{2}=4r^{2}-r^{2}=3r^{2}
\]
代入 \(AE=2\sqrt3\)，得 \(r=2\)，从而 \(AO=4\).

弦 \(MN=2\sqrt3\)，且圆心到弦的垂线为 \(OD\)，所以
\[
OD=\sqrt{r^{2}-\left(\frac{MN}{2}\right)^{2}}=\sqrt{4-3}=1
\]
因此 \(AD=AO+OD=5\). 在直角三角形 \(AOE\) 中，\(\sin\angle OAE=OE/AO=1/2\)，故 \(\angle BAD=30^{\circ}\). 于是
\(AB=\frac{AD}{\cos30^{\circ}}=\frac{10\sqrt3}{3}\)，\(BC=2BD=2AD\tan30^{\circ}=\frac{10\sqrt3}{3}\).
由 \(EF\parallel BC\)，三角形 \(AEF\) 与三角形 \(ABC\) 相似，且相似比为
\[
\frac{AE}{AB}=\frac{2\sqrt3}{10\sqrt3/3}=\frac35
\]
所以
\[
S_{\triangle ABC}=\frac12\cdot\frac{10\sqrt3}{3}\cdot5=\frac{25\sqrt3}{3}
\]
\[
S_{\triangle AEF}=\left(\frac35\right)^{2}S_{\triangle ABC}=3\sqrt3
\]
故
\[
S_{EBCF}=S_{\triangle ABC}-S_{\triangle AEF}=\frac{16\sqrt3}{3}
\]
\end{enumerate}
\end{solution}

\begin{problem}
在直角坐标系 \(xOy\) 中，曲线 \(C_1: \begin{cases}
x = t\cos \alpha,\\
y = t\sin \alpha,
\end{cases}\)（\(t\) 为参数，\(t \neq 0\)），其中 \(0 \leqslant \alpha < \pi\). 在以 \(O\) 为极点，\(x\) 轴正半轴为极轴的极坐标系中，曲线 \(C_2: \rho = 2\sin \theta\)，\(C_3: \rho = 2\sqrt{3}\cos \theta\).
\begin{enumerate}
\item 求 \(C_2\) 与 \(C_3\) 交点的直角坐标；
\item 若 \(C_1\) 与 \(C_2\) 相交于点 \(A\)，\(C_1\) 与 \(C_3\) 相交于点 \(B\)，求 \(|AB|\) 的最大值.
\end{enumerate}
\end{problem}

\begin{solution}
\begin{enumerate}
\item 由极坐标与直角坐标的关系，\(\rho\sin\theta=y\)，\(\rho\cos\theta=x\). 因而
\(C_{2}:x^{2}+y^{2}=2y\)，\(C_{3}:x^{2}+y^{2}=2\sqrt3x\).
两式相减得 \(y=\sqrt3x\). 代回第一式，得
\[
4x^{2}=2\sqrt3x
\]
所以 \(x=0\) 或 \(x=\dfrac{\sqrt3}{2}\). 对应的交点为
\((0,0)\)，\(\left(\frac{\sqrt3}{2},\frac32\right)\).

\item 由参数方程，\(C_{1}\) 上的点可写为 \((t\cos\alpha,t\sin\alpha)\). 将其代入 \(C_{2}\) 的直角坐标方程，且 \(t\ne0\)，得
\[
t^{2}=2t\sin\alpha\Longrightarrow t_{A}=2\sin\alpha
\]
同理，将其代入 \(C_{3}\)，得
\[
t^{2}=2\sqrt3t\cos\alpha\Longrightarrow t_{B}=2\sqrt3\cos\alpha
\]
两点在同一直线 \(C_{1}\) 上，故
\[
|AB|=|t_{A}-t_{B}|=2|\sin\alpha-\sqrt3\cos\alpha|
\]
由柯西不等式，
\[
|\sin\alpha-\sqrt3\cos\alpha|\leqslant\sqrt{1+3}=2
\]
等号可在 \(\alpha=\dfrac{5\pi}{6}\) 时取得，所以 \(|AB|\) 的最大值为 \(4\).
\end{enumerate}
\end{solution}

\begin{problem}
设 \(a\)，\(b\)，\(c\)，\(d\) 均为正数，且 \(a + b = c + d\)，证明：
\begin{enumerate}
\item 若 \(ab > cd\)，则 \(\sqrt{a} + \sqrt{b} > \sqrt{c} + \sqrt{d}\)；
\item \(\sqrt{a} + \sqrt{b} > \sqrt{c} + \sqrt{d}\) 是 \(|a - b| < |c - d|\) 的充要条件.
\end{enumerate}
\end{problem}

\begin{solution}
\begin{enumerate}
\item 因为四个数均为正数，可以对两边平方. 又由 \(a+b=c+d\)，有
\((\sqrt a+\sqrt b)^{2}=a+b+2\sqrt{ab}\)，\((\sqrt c+\sqrt d)^{2}=c+d+2\sqrt{cd}\).
条件 \(ab>cd\) 给出 \(\sqrt{ab}>\sqrt{cd}\)，从而
\[
(\sqrt a+\sqrt b)^{2}>(\sqrt c+\sqrt d)^{2}
\]
两边均为正数，所以 \(\sqrt a+\sqrt b>\sqrt c+\sqrt d\).

\item 仍由 \(a+b=c+d=:s\)，有
\((a-b)^{2}=s^{2}-4ab\)，\((c-d)^{2}=s^{2}-4cd\).
因此
\[
|a-b|<|c-d|\Longleftrightarrow ab>cd
\]
另一方面，由上面的平方展开式，
\[
\sqrt a+\sqrt b>\sqrt c+\sqrt d\Longleftrightarrow ab>cd
\]
所以
\[
\sqrt a+\sqrt b>\sqrt c+\sqrt d
\Longleftrightarrow |a-b|<|c-d|
\]
\end{enumerate}
\end{solution}
